\documentclass[aspectratio=169]{beamer}
% talk-common.tex —— 各场分享共用的导言区配置
% 放在仓库根目录，由各个入口 .tex 通过 % talk-common.tex —— 各场分享共用的导言区配置
% 放在仓库根目录，由各个入口 .tex 通过 % talk-common.tex —— 各场分享共用的导言区配置
% 放在仓库根目录，由各个入口 .tex 通过 \input{talk-common.tex} 引用。
% 编译时需要把仓库根目录加入 TEXINPUTS，详见根目录 AGENTS.md 第 2 节。

\usetheme[customfont,showmaxslides,nofooterlogo]{pureminimalistic}
\usefonttheme{professionalfonts}
\usepackage[UTF8,fontset=none]{ctex}
\usepackage{graphicx}
\usepackage{listings}
\usepackage{appendixnumberbeamer}

\setsansfont{Fira Sans}
\setmonofont{Fira Mono}
\setCJKmainfont{Noto Sans CJK SC}
\setCJKsansfont{Noto Sans CJK SC}
\setCJKmonofont{Noto Sans CJK SC}

% 默认不引用上游示例 logo，有真实素材时再统一配置
\renewcommand{\logotitle}{}
\renewcommand{\logoheader}{}
\renewcommand{\logofooter}{}
\renewcommand{\pageword}{}

% 字号设置必须放在加载主题之后，避免被主题覆盖
\setbeamerfont{title}{size=\LARGE,shape=\upshape,series=\bfseries}
\setbeamerfont{presentation title}{shape=\upshape}
\setbeamerfont{frametitle}{size=\Large,shape=\upshape,series=\bfseries}

% 对浅色/深色都保持明确的代码对比度
\definecolor{TalkCodeBg}{HTML}{F3F4F6}
\definecolor{TalkCodeFg}{HTML}{17212B}
\definecolor{TalkCodeKeyword}{HTML}{174EA6}
\definecolor{TalkCodeComment}{HTML}{47634D}
\definecolor{TalkCodeString}{HTML}{8A3014}
\lstdefinestyle{talkcode}{
  basicstyle=\ttfamily\small\color{TalkCodeFg},
  backgroundcolor=\color{TalkCodeBg},
  keywordstyle=\bfseries\color{TalkCodeKeyword},
  commentstyle=\color{TalkCodeComment},
  stringstyle=\color{TalkCodeString},
  columns=fullflexible,
  keepspaces=true,
  showstringspaces=false,
  breaklines=true,
  tabsize=4,
  numbers=none,
  frame=none
}
\lstset{style=talkcode}
 引用。
% 编译时需要把仓库根目录加入 TEXINPUTS，详见根目录 AGENTS.md 第 2 节。

\usetheme[customfont,showmaxslides,nofooterlogo]{pureminimalistic}
\usefonttheme{professionalfonts}
\usepackage[UTF8,fontset=none]{ctex}
\usepackage{graphicx}
\usepackage{listings}
\usepackage{appendixnumberbeamer}

\setsansfont{Fira Sans}
\setmonofont{Fira Mono}
\setCJKmainfont{Noto Sans CJK SC}
\setCJKsansfont{Noto Sans CJK SC}
\setCJKmonofont{Noto Sans CJK SC}

% 默认不引用上游示例 logo，有真实素材时再统一配置
\renewcommand{\logotitle}{}
\renewcommand{\logoheader}{}
\renewcommand{\logofooter}{}
\renewcommand{\pageword}{}

% 字号设置必须放在加载主题之后，避免被主题覆盖
\setbeamerfont{title}{size=\LARGE,shape=\upshape,series=\bfseries}
\setbeamerfont{presentation title}{shape=\upshape}
\setbeamerfont{frametitle}{size=\Large,shape=\upshape,series=\bfseries}

% 对浅色/深色都保持明确的代码对比度
\definecolor{TalkCodeBg}{HTML}{F3F4F6}
\definecolor{TalkCodeFg}{HTML}{17212B}
\definecolor{TalkCodeKeyword}{HTML}{174EA6}
\definecolor{TalkCodeComment}{HTML}{47634D}
\definecolor{TalkCodeString}{HTML}{8A3014}
\lstdefinestyle{talkcode}{
  basicstyle=\ttfamily\small\color{TalkCodeFg},
  backgroundcolor=\color{TalkCodeBg},
  keywordstyle=\bfseries\color{TalkCodeKeyword},
  commentstyle=\color{TalkCodeComment},
  stringstyle=\color{TalkCodeString},
  columns=fullflexible,
  keepspaces=true,
  showstringspaces=false,
  breaklines=true,
  tabsize=4,
  numbers=none,
  frame=none
}
\lstset{style=talkcode}
 引用。
% 编译时需要把仓库根目录加入 TEXINPUTS，详见根目录 AGENTS.md 第 2 节。

\usetheme[customfont,showmaxslides,nofooterlogo]{pureminimalistic}
\usefonttheme{professionalfonts}
\usepackage[UTF8,fontset=none]{ctex}
\usepackage{graphicx}
\usepackage{listings}
\usepackage{appendixnumberbeamer}

\setsansfont{Fira Sans}
\setmonofont{Fira Mono}
\setCJKmainfont{Noto Sans CJK SC}
\setCJKsansfont{Noto Sans CJK SC}
\setCJKmonofont{Noto Sans CJK SC}

% 默认不引用上游示例 logo，有真实素材时再统一配置
\renewcommand{\logotitle}{}
\renewcommand{\logoheader}{}
\renewcommand{\logofooter}{}
\renewcommand{\pageword}{}

% 字号设置必须放在加载主题之后，避免被主题覆盖
\setbeamerfont{title}{size=\LARGE,shape=\upshape,series=\bfseries}
\setbeamerfont{presentation title}{shape=\upshape}
\setbeamerfont{frametitle}{size=\Large,shape=\upshape,series=\bfseries}

% 对浅色/深色都保持明确的代码对比度
\definecolor{TalkCodeBg}{HTML}{F3F4F6}
\definecolor{TalkCodeFg}{HTML}{17212B}
\definecolor{TalkCodeKeyword}{HTML}{174EA6}
\definecolor{TalkCodeComment}{HTML}{47634D}
\definecolor{TalkCodeString}{HTML}{8A3014}
\lstdefinestyle{talkcode}{
  basicstyle=\ttfamily\small\color{TalkCodeFg},
  backgroundcolor=\color{TalkCodeBg},
  keywordstyle=\bfseries\color{TalkCodeKeyword},
  commentstyle=\color{TalkCodeComment},
  stringstyle=\color{TalkCodeString},
  columns=fullflexible,
  keepspaces=true,
  showstringspaces=false,
  breaklines=true,
  tabsize=4,
  numbers=none,
  frame=none
}
\lstset{style=talkcode}


\title[技术部第五次分享]{技术部第五次技术分享会}
\subtitle{枚举、泛型与常用标准库}
\author{姚雨森}
\institute{技术部 分管主席}
\date{}

\begin{document}

\maketitle

\section{枚举：表示有限的一组选择}

\begin{frame}[fragile]{用整数表示状态，需要记住编号的含义}
  一个订单可能处于待付款、已付款、已发货或已取消状态。
  \begin{lstlisting}[language=Java]
int status = 0;

if (status == 1) {
    System.out.println("Ready to ship");
}

status = 7;
  \end{lstlisting}
  \begin{itemize}
    \item 阅读代码时，需要另外查找：\texttt{0} 和 \texttt{1} 分别代表什么？
    \item \texttt{int} 允许许多整数，但业务只需要四种状态。
    \item 即使 \texttt{7} 没有业务含义，赋值语句仍能通过类型检查。
  \end{itemize}
\end{frame}

\begin{frame}[fragile]{换成字符串更直观，但仍然可能写错}
  \begin{lstlisting}[language=Java]
String status = "PENDING";
status = "PAID";

status = "PAIID";
status = "paid";
  \end{lstlisting}
  \begin{itemize}
    \item 名称比编号容易理解，但任何字符串都可以赋给 \texttt{status}。
    \item 拼写错误、大小写不一致，都可能让后面的判断失效。
    \item 我们希望同时表达\alert{每个选项的含义}和\alert{允许使用的选项}。
  \end{itemize}
\end{frame}

\begin{frame}{枚举：为有限的一组选择定义类型}
  订单状态只需要以下四个选项：
  \par\bigskip
  \begin{center}
    \renewcommand{\arraystretch}{1.35}
    \begin{tabular}{ll}
      \textbf{含义} & \textbf{枚举常量名称} \\
      \hline
      待付款 & \texttt{PENDING} \\
      已付款 & \texttt{PAID} \\
      已发货 & \texttt{SHIPPED} \\
      已取消 & \texttt{CANCELLED} \\
    \end{tabular}
  \end{center}
  \bigskip
  \textbf{枚举（enum）}将这些选项列在一起，定义一个专门的类型。
  \par\medskip
  星期、方向、任务状态等有限选项，也可以用枚举表达。
\end{frame}

\begin{frame}[fragile]{使用 enum 定义订单状态类型}
  将下面的定义放在 \texttt{OrderStatus.java} 中：
  \begin{lstlisting}[language=Java]
public enum OrderStatus {
    PENDING,
    PAID,
    SHIPPED,
    CANCELLED
}
  \end{lstlisting}
  \begin{itemize}
    \item \texttt{enum} 用于声明枚举；\texttt{OrderStatus} 是新类型的名称。
    \item 花括号中列出枚举常量，使用逗号分隔。
    \item 类型名称通常使用大驼峰；常量通常使用大写字母和下划线。
    \item 使用常量时写成\alert{类型名.常量名}，例如 \texttt{OrderStatus.PAID}。
  \end{itemize}
\end{frame}

\begin{frame}[fragile]{像其他类型一样声明变量和赋值}
  以下语句写在 \texttt{Main.java} 的 \texttt{main} 方法中，
  两个源文件放在同一目录。
  \begin{lstlisting}[language=Java]
OrderStatus status = OrderStatus.PENDING;
System.out.println(status); // PENDING

status = OrderStatus.PAID;
System.out.println(status); // PAID
  \end{lstlisting}
  \begin{itemize}
    \item 左边的 \texttt{OrderStatus} 指定变量类型。
    \item 右边的 \texttt{OrderStatus.PENDING} 是该类型的一个常量。
    \item 变量可以重新赋值；这里是改为保存另一个枚举常量。
  \end{itemize}
\end{frame}

\begin{frame}[fragile]{类型检查帮助我们排除无效的赋值}
  下面的错误写法用于说明类型限制，\alert{不能通过编译}：
  \begin{lstlisting}[language=Java]
OrderStatus status = OrderStatus.PENDING;

status = 1;                  // Wrong type
status = "PAID";             // Wrong type
status = OrderStatus.PAIID;   // No such constant
  \end{lstlisting}
  \begin{itemize}
    \item 枚举常量有自己的类型；字符串 \texttt{"PAID"} 并不是它。
    \item 拼错常量名称时，编译器就能指出问题。
    \item 另一种枚举类型的常量，即使同名，也不能赋给这个变量。
  \end{itemize}
  \medskip
  补充：枚举变量也可以为 \texttt{null}，表示没有引用枚举常量；
  本例始终给它赋明确的状态。
\end{frame}

\begin{frame}[fragile]{使用 == 判断是否为某个枚举常量}
  \begin{lstlisting}[language=Java]
OrderStatus status = OrderStatus.PAID;

if (status == OrderStatus.PAID) {
    System.out.println("Ready to ship");
}

if (status != OrderStatus.CANCELLED) {
    System.out.println("Order is not cancelled");
}
  \end{lstlisting}
  \begin{itemize}
    \item \texttt{==} 判断是否为同一个枚举常量；\texttt{!=} 判断是否不同。
    \item 直接比较枚举值，不需要先转换成编号或字符串。
    \item 字符串内容比较仍使用 \texttt{equals}，不要把这里的规则照搬过去。
  \end{itemize}
\end{frame}

\begin{frame}[fragile]{枚举作为参数：明确方法接收哪一种状态}
  假设本例规定：只有待付款的订单可以取消。
  \par\medskip
  在 \texttt{Main} 类中定义方法：
  \begin{lstlisting}[language=Java]
static boolean canCancel(OrderStatus status) {
    return status == OrderStatus.PENDING;
}
  \end{lstlisting}
  在 \texttt{main} 方法中调用：
  \begin{lstlisting}[language=Java]
boolean first = canCancel(OrderStatus.PENDING); // true
boolean second = canCancel(OrderStatus.SHIPPED); // false
  \end{lstlisting}
  \medskip
  参数的类型说明需要什么数据；方法体决定这组数据对应的业务规则。
\end{frame}

% \begin{frame}[fragile]{枚举作为返回值：让调用方继续判断}
%   为演示返回值，假设新订单统一从待付款状态开始。
%   \par\medskip
%   在 \texttt{Main} 类中定义方法：
%   \begin{lstlisting}[language=Java]
% static OrderStatus initialStatus() {
%     return OrderStatus.PENDING;
% }
%   \end{lstlisting}
%   在 \texttt{main} 方法中使用返回值：
%   \begin{lstlisting}[language=Java]
% OrderStatus status = initialStatus();
% if (canCancel(status)) {
%     System.out.println("Cancellation is allowed");
% }
%   \end{lstlisting}
%   \medskip
%   返回类型写 \texttt{OrderStatus}，调用方就可以用同类型变量接收结果。
% \end{frame}

% \begin{frame}{枚举限定可选值，业务规则仍由程序实现}
%   \begin{itemize}
%     \item \textbf{定义类型：}用 \texttt{enum} 列出有限的一组常量。
%     \item \textbf{保存与比较：}用枚举变量保存选项，用 \texttt{==} 判断常量。
%     \item \textbf{连接方法：}枚举可以作为参数，也可以作为返回值。
%   \end{itemize}
%   \bigskip
%   \begin{block}{思考：下面这次赋值能通过编译吗？}
%     已发货的订单变量重新赋为 \texttt{OrderStatus.PENDING}。
%   \end{block}
%   \medskip
%   能：两个常量属于同一种类型。
%   枚举限定有哪些状态，\alert{不会自动限制状态之间如何转换}。
% \end{frame}

\section{错误处理：报告失败并由调用方处理}

\begin{frame}[fragile]{商品价格需要满足非负的要求}
  回顾：通过 setter 修改商品的价格。本例价格用整数表示，单位为分。
  \begin{lstlisting}[language=Java]
public class Product {
    private int price;

    public void setPrice(int price) {
        this.price = price;
    }
}
  \end{lstlisting}
  \begin{itemize}
    \item 字段设为 \texttt{private}，类外通过 \texttt{setPrice} 修改它。
    \item 目前调用 \texttt{setPrice(-10)}，仍然会把价格改为负数。
    \item 我们希望\alert{拒绝非法价格，并让调用方知道操作失败}。
  \end{itemize}
\end{frame}

\begin{frame}[fragile]{使用 throw 抛出异常，拒绝非法参数}
  将 \texttt{Product} 中的 setter 改为：
  \begin{lstlisting}[language=Java]
public void setPrice(int price) {
    if (price < 0) {
        throw new IllegalArgumentException(
            "Price must be non-negative");
    }
    this.price = price;
}
  \end{lstlisting}
  \begin{itemize}
    \item \texttt{IllegalArgumentException} 表示参数不符合方法的要求。
    \item \texttt{new} 创建异常对象；字符串说明失败的原因。
    \item \texttt{throw} 抛出异常，本次调用不能正常完成。
    \item 检查在赋值之前：负数不会写入字段，\texttt{0} 则是合法价格。
  \end{itemize}
\end{frame}

% \begin{frame}{抛出异常后，不会继续执行后面的赋值}
%   \begin{columns}[T,onlytextwidth]
%     \begin{column}{0.47\textwidth}
%       \begin{block}{传入 1000：满足要求}
%         \begin{enumerate}
%           \item 判断价格是否小于 \texttt{0}：否。
%           \item 将 \texttt{1000} 写入字段。
%           \item 方法正常结束，回到调用方。
%         \end{enumerate}
%       \end{block}
%     \end{column}
%     \begin{column}{0.49\textwidth}
%       \begin{block}{传入 -10：不满足要求}
%         \begin{enumerate}
%           \item 判断价格是否小于 \texttt{0}：是。
%           \item 创建并抛出异常。
%           \item 跳过赋值，异常向调用方传播。
%         \end{enumerate}
%       \end{block}
%     \end{column}
%   \end{columns}
%   \bigskip
%   \alert{异常表示这次操作失败}，不是正常返回的结果。
%   \par\medskip
%   当前 setter 在检查失败时没有修改字段，原价格得以保留。
% \end{frame}

\begin{frame}{异常类型帮助我们识别失败原因}
  \begin{center}
    \renewcommand{\arraystretch}{1.4}
    \begin{tabular}{ll}
      \textbf{异常类型} & \textbf{典型原因} \\
      \hline
      \texttt{IllegalArgumentException} & 参数不符合要求 \\
      \texttt{NumberFormatException} & 文本不能转换为要求的数值 \\
      \texttt{NullPointerException} & 对 \texttt{null} 进行需要对象的操作 \\
      \texttt{IndexOutOfBoundsException} & 索引超出允许范围 \\
    \end{tabular}
  \end{center}
  \bigskip
  \begin{itemize}
    \item 参数校验失败时，可以主动抛出合适的异常。
    \item 调用库方法或执行其他操作时，也可能出现异常。
    \item 类型说明问题的类别，异常消息提供本次失败的具体信息。
  \end{itemize}
\end{frame}

\begin{frame}[fragile]{异常没有被处理，调用方也不能照常继续}
  在 \texttt{main} 方法中调用修改后的 setter：
  \begin{lstlisting}[language=Java]
Product product = new Product();
product.setPrice(1000);
product.setPrice(-10);
System.out.println("Price updated");
  \end{lstlisting}
  \begin{itemize}
    \item 第一次设置成功；第二次在 setter 中抛出异常。
    \item 异常传播到 \texttt{main}，最后的输出语句不会执行。
    \item 始终没有匹配的处理时，当前线程终止；本例程序随之结束。
  \end{itemize}
  \medskip
  控制台会给出异常类型、消息和调用堆栈；
  可以从中找到自己代码的文件名与行号。
\end{frame}

\begin{frame}[fragile]{在调用方使用 try-catch 处理设置失败}
  以下代码写在 \texttt{main} 方法中：
  \begin{lstlisting}[language=Java]
Product product = new Product();
product.setPrice(1000);
try {
    product.setPrice(-10);
    System.out.println("Price updated");
} catch (IllegalArgumentException e) {
    System.out.println(e.getMessage());
}
System.out.println("Continue");
  \end{lstlisting}
  \begin{itemize}
    \item \texttt{try} 中放置可能失败的操作。
    \item \texttt{catch} 捕获匹配的异常，变量 \texttt{e} 接收异常对象。
    \item \texttt{e.getMessage()} 取得异常消息，此处用它给出失败提示。
  \end{itemize}
\end{frame}

\begin{frame}{对照两条路径：正常执行与异常处理}
  将上一页的 \texttt{-10} 分别改为 \texttt{2000} 和 \texttt{-10}：
  \par\bigskip
  \begin{columns}[T,onlytextwidth]
    \begin{column}{0.47\textwidth}
      \begin{block}{传入 2000：正常完成}
        \begin{enumerate}
          \item 设置价格成功。
          \item 输出 \texttt{Price updated}。
          \item 跳过 \texttt{catch}。
          \item 输出 \texttt{Continue}。
        \end{enumerate}
      \end{block}
    \end{column}
    \begin{column}{0.49\textwidth}
      \begin{block}{传入 -10：进入异常处理}
        \begin{enumerate}
          \item setter 抛出异常。
          \item 跳过 \texttt{try} 中剩余语句。
          \item 进入匹配的 \texttt{catch}，输出原因。
          \item 输出 \texttt{Continue}。
        \end{enumerate}
      \end{block}
    \end{column}
  \end{columns}
  \bigskip
  本例处理后，继续执行整个 \texttt{try-catch} 后面的语句，
  \alert{不会回到出错位置接着执行}。
\end{frame}

\begin{frame}{setter 检查规则，调用方处理失败}
  \begin{itemize}
    \item \textbf{setter：}检查价格，拒绝负数，并抛出异常报告失败。
    \item \textbf{调用方：}根据场景提示用户、重新输入，或取消这次操作。
    \item \textbf{自己的代码写错时：}修正非法参数的来源，避免反复触发同一问题。
  \end{itemize}
  \bigskip
  \begin{block}{捕获异常不会自动修复数据，也不会自动重试}
    本例价格仍为原来的 \texttt{1000}，是因为 setter 先检查、后赋值。
    异常机制本身不会撤销此前已经完成的修改。
  \end{block}
  \medskip
  不要写空的 \texttt{catch} 来隐藏失败；在能作出合理处理的位置捕获异常。
\end{frame}

\begin{frame}{补充：有些异常要求捕获或显式声明}
  \begin{itemize}
    \item Java 对部分异常有编译要求：调用方必须捕获，
      或在方法上用 \texttt{throws} 声明可能向外传播。
    \item 本例的 \texttt{IllegalArgumentException} 没有这项强制要求，
      也可以使用 \texttt{try-catch} 处理。
  \end{itemize}
  \bigskip
  \begin{block}{区分两个关键字}
    \texttt{throw}：实际抛出异常。\par\medskip
    \texttt{throws}：写在方法声明中，说明可能向外传播的异常，
    不代表已经处理了异常。
  \end{block}
  \medskip
  课后了解：受检异常与非受检异常的分类，以及 \texttt{throws} 的使用。
\end{frame}

\end{document}
